% Resultado de máxima entropía recuperado del propietario de incidencia.
% Se desarrolla aquí su prueba variacional para hacerla local al artículo.
% No se atribuye novedad al resultado ni a su formulación convencional.
\subsection{Distribution stationnaire et caractérisation variationnelle}
\label{v:subsec:parry-explicita}

Soit \(X_{\rm av}\) l’espace des suites bilatérales sur
\(\{\mathrm{ar},\mathrm{vol}\}\) dont les paires consécutives ont une
incidence positive dans \(F_{\rm av}\). C’est l’enveloppe à mémoire un
du calendrier des modes : elle conserve ses paires admissibles, tout en permettant
toutes les concaténations compatibles avec ces paires. Sa mesure
stationnaire pondère cet espace d’histoires ; la fréquence chronologique
de l’orbite \((\Sigma,\Pi,\Pi)\) reste, séparément,
\((1/3,2/3)\).

\begin{proposition}[Caractérisation explicite de la mesure d’entropie maximale]
\label{v:prop:parry-explicita}
Avec \(r=(1,\varphi)\), la matrice stochastique et sa distribution
stationnaire sont
\begin{equation}
 P_{ij}=\frac{(F_{\rm av})_{ij}r_j}{\varphi r_i},\qquad
 P=\begin{pmatrix}0&1\\\varphi^{-2}&\varphi^{-1}\end{pmatrix},
 \qquad
 \pi_{\rm av}=\frac{(1,\varphi^2)}{1+\varphi^2}.
 \label{v:eq:parry-transicion-estacionaria}
\end{equation}
La mesure stationnaire de Markov \(\mu\) déterminée par ces données
a, pour chaque cylindre admissible de \(n\) transitions,
\begin{equation}
 \mu[i_0\ldots i_n]
 =\frac{r_{i_0}r_{i_n}}{(1+\varphi^2)\varphi^n}.
 \label{v:eq:parry-cilindros}
\end{equation}
C’est l’unique mesure stationnaire sur \(X_{\rm av}\) ayant un taux
d’entropie maximal, et ce taux est \(\log\varphi\).
\end{proposition}

\begin{proof}
L’identité \(\varphi^2=\varphi+1\) implique
\(\varphi^{-2}+\varphi^{-1}=1\), de sorte que les lignes de \(P\)
ont pour somme un. La substitution directe vérifie
\(\pi_{\rm av}P=\pi_{\rm av}\). Les deux transitions croisées
ont une probabilité positive ; la chaîne est irréductible, et le maintien
en \(\mathrm{vol}\) la rend apériodique. La distribution stationnaire
s’obtient aussi en résolvant les deux équations linéaires et la
normalisation, qui donnent l’unique solution imprimée. La multiplication de
\(\pi_{i_0}\prod_{k=0}^{n-1}P_{i_k i_{k+1}}\) annule les facteurs
intermédiaires \(r_{i_k}\) et démontre \eqref{v:eq:parry-cilindros}.

Pour la maximalité, soit \(\nu\) une mesure stationnaire quelconque
sur \(X_{\rm av}\) et soit
\(q(j\mid\mathcal P)\) la distribution conditionnelle de \(X_1\)
sachant le passé complet \(\mathcal P=(\ldots,X_{-1},X_0)\).
Son support est contenu dans celui de la ligne \(P_{X_0,\cdot}\).
Pour des distributions finies, la divergence relative
\(D(q\Vert p)=\sum_jq_j\log(q_j/p_j)\) est non négative, et
s’annule exactement lorsque \(q=p\). Cette propriété s’obtient par
la concavité du logarithme, avec la convention \(0\log0=0\).
Sur les paires admissibles,
\[
 \log P_{ij}=\log r_j-\log r_i-\log\varphi.
\]
La stationnarité annule l’espérance de
\(\log r_{X_1}-\log r_{X_0}\). Par conséquent,
\begin{equation}
 h(\nu)=H_\nu(X_1\mid\mathcal P)
 =\log\varphi-
 \mathbb E_\nu D\!\left(q(\cdot\mid\mathcal P)
                  \Vert P_{X_0,\cdot}\right)
 \leq\log\varphi.
 \label{v:eq:parry-deficit-entropia}
\end{equation}
L’identité entre taux et entropie conditionnelle du passé s’obtient
au moyen de la règle de la chaîne sur les blocs finis et de la limite
décroissante de leurs entropies conditionnelles ; l’alphabet fini garantit
qu’elles sont toutes bornées entre zéro et \(\log2\).

La mesure \(\mu\) construite est markovienne avec transition \(P\),
de sorte que le terme de divergence s’annule et
\(h(\mu)=\log\varphi\). Si une autre mesure stationnaire atteint
ce taux, la divergence s’annule presque sûrement : sa distribution
conditionnelle du futur dépend seulement de \(X_0\) et coïncide avec \(P\).
C’est donc la même chaîne stationnaire, dont l’unique distribution initiale
est \(\pi_{\rm av}\). Cela démontre l’unicité.
\end{proof}

La formule \eqref{v:eq:parry-deficit-entropia} précise ce qui est maximisé :
l’entropie par transition de l’enveloppe symbolique déclarée. Le rapport
\(\mu_{\rm ar}/\mu_{\rm vol}=\varphi^{-2}\) provient de ses poids
stationnaires. L’orbite chronologique périodique a un taux d’entropie
nul et maintient sa propre répartition temporelle ; les deux objets restent
définis avec leurs domaines respectifs.
